\documentclass[11pt,a4paper]{article}
% Estilo comum dos PDFs individuais da Sprint 1.
\usepackage[utf8]{inputenc}
\usepackage[T1]{fontenc}
\usepackage[brazilian]{babel}
\usepackage[a4paper,margin=2.35cm]{geometry}
\usepackage{xcolor}
\usepackage{hyperref}
\usepackage{enumitem}
\usepackage{listings}
\usepackage{parskip}
\usepackage{needspace}
\definecolor{azul}{RGB}{20,63,112}
\hypersetup{colorlinks=true,linkcolor=azul,urlcolor=azul}
\setlist[itemize]{leftmargin=*,itemsep=3pt}
\lstset{basicstyle=\ttfamily\small,frame=single,breaklines=true,columns=fullflexible}
\sloppy
\newcommand{\codigo}[1]{\texttt{#1}}
\newcommand{\titulo}[2]{\begin{center}{\LARGE\bfseries #1\par}\vspace{4pt}{\large #2\par}\vspace{8pt}\rule{0.8\textwidth}{0.4pt}\par\vspace{5pt}{\small Matheus Moreli Durão Silva -- Matrícula 12621CBS217\par Cibersegurança - Bacharelado - Noturno\par}\end{center}}
\newcommand{\cardfonte}[1]{\par\smallskip\noindent\textit{Referência: }\href{#1}{Trello - enunciado e checklist do card}.}

\begin{document}\titulo{S1\_06 - Entrada e saída}{Conversar com quem usa o programa}
\section*{Entregável correto}
O programa deve solicitar ID, nome, responsável e localização de um ativo e apresentar esses dados de maneira formatada. O foco é compreender que a entrada é texto até que seja convertida e que a saída deve ser legível para a pessoa usuária.
\section*{Passo a passo}
\codigo{input("Pergunta: ")} mostra uma pergunta e espera o teclado. O valor devolvido é uma \codigo{str}; por isso, quando ID representa número, usa-se \codigo{int(...)}. Uma f-string começa com \codigo{f} antes das aspas e coloca expressões entre chaves. Ela evita concatenar muitos trechos e preserva a leitura do texto.
\begin{lstlisting}
id_ativo = int(input("ID: "))
nome = input("Nome: ")
responsavel = input("Responsavel: ")
localizacao = input("Localizacao: ")
print(f"Ativo {id_ativo}: {nome}\nResponsavel: {responsavel}\nLocal: {localizacao}")
\end{lstlisting}
O \codigo{\textbackslash n} cria uma nova linha. Se o usuário digitar letras para o ID, a conversão falha; na S1\_15 essa situação passa a ser tratada com \codigo{try/except}.
\section*{Uso no projeto final}
Cadastro, busca e atualização no \codigo{app.py} usam a mesma sequência: pedir, converter quando necessário, validar e mostrar uma resposta. A diferença é que o projeto armazena o resultado para uso posterior.
\section*{Checklist demonstrado}
Entrada com \codigo{input}, saída com \codigo{print}, conversão de tipos, f-strings e interação usuário-sistema. Teste com valores reais e confira que todos os quatro campos aparecem, sem substituir o ID numérico por texto.
\cardfonte{https://trello.com/c/fFlJ9zCj}
\end{document}
